% Generated draft: requires mathematical and source review.
\documentclass[11pt]{article}
\usepackage[margin=25mm]{geometry}
\usepackage{fontspec}
\setmainfont{Latin Modern Roman}
\usepackage{amsmath,amssymb,mathtools,unicode-math}
\setmathfont{Latin Modern Math}
\usepackage{enumitem,needspace,xurl,hyperref}
\setlength{\parindent}{0pt}
\setlength{\parskip}{5pt}
\setlength{\emergencystretch}{3em}
\newcommand{\N}{\mathbb{N}}
\newcommand{\Z}{\mathbb{Z}}
\newcommand{\Q}{\mathbb{Q}}
\newcommand{\R}{\mathbb{R}}
\newcommand{\C}{\mathbb{C}}
\newcommand{\F}{\mathbb{F}}
\newcommand{\A}{\mathbb{A}}
\newcommand{\PP}{\mathbb{P}}
\newcommand{\E}{\mathbb{E}}
\newcommand{\eps}{\varepsilon}
\newcommand{\dd}{\,\mathrm{d}}
\begin{document}
\section*{UnsolvedMath ID: 30002195}
\textbf{Source problem code:} catalogue entry
\subsection{Definitions and mathematical statement}
A generalized Cartan matrix indexed by a finite set \(I\) is an integer matrix \(C=(c_{ij})_{i,j\in I}\) such that \(c_{ii}=2\), \(c_{ij}\leq 0\) for \(i\ne j\), and \(c_{ij}=0\) if and only if \(c_{ji}=0\). A Cartan scheme consists of such an index set, a nonempty object set \(A\), involutions \(\rho_i:A\to A\) for \(i\in I\), and generalized Cartan matrices \(C^a\) for \(a\in A\), satisfying
\[
 c^a_{ij}=c^{\rho_i(a)}_{ij}\qquad(a\in A,\ i,j\in I).
\]
The supplied record does not specify the additional hypotheses or the meaning of finiteness used in the referenced classification; finiteness of the object set alone must not be substituted for those hypotheses.

Write \(\mathcal C\) schematically for the objects in the referenced classification of finite Cartan schemes, and \(\mathcal A\subseteq\mathcal C\) for those that occur in the referenced classification of arithmetic root systems, using the intended correspondence and equivalence conventions. These classifications, conventions, and correspondence are not specified in the supplied record, so this notation does not independently define an explicit class. The intended exceptional class is
\[
\mathcal E=\mathcal C\setminus\mathcal A.
\]
A Lie algebra over a field \(k\) of positive characteristic is a \(k\)-vector space \(L\) with a bilinear alternating bracket satisfying the Jacobi identity, where \(\operatorname{char}(k)=p\) for a prime \(p\).

\textbf{Source question.} Are the objects in \(\mathcal E\) related to Lie algebras over fields of positive characteristic?

The word ``related'' is not given a mathematical definition in the supplied record. In particular, no construction, correspondence, invariant, characteristic, dimension condition, or quantifier specifying which exceptional objects should have such a relation is supplied. Consequently, the question above preserves the corrected source formulation but is not yet a fully specified conjecture.

The earlier wording concerns ``remaining finite-dimensional Nichols algebras of diagonal type'' and unidentified Cartan-scheme objects. The supplied verification explicitly corrects that target to the additional Cartan-scheme objects absent from the arithmetic-root-system classification. Thus the earlier wording supplies context, but does not justify identifying every exceptional object with a finite-dimensional Nichols algebra.
\subsection{Short English statement}
The question asks whether the extra objects in the finite Cartan-scheme classification, beyond those occurring in the arithmetic-root-system classification, have a connection with Lie algebras in positive characteristic. The classifications, the intended connection, and its quantifiers require clarification before this becomes a precise conjecture.
\subsection{Sources}
\begin{itemize}
\item Supplied local registry record, permanent UnsolvedMath ID \(30002195\), title ``Cartan Schemes and Positive-Characteristic Lie Algebras,'' source problem code \texttt{catalogue entry}, in the Oberwolfach Reports Open Problems collection. No individual proposer is identified; the supplied proposed year is 2012.
\item Supplied report citation: \emph{Mini-Workshop: Nichols Algebras and Weyl Groupoids (2013)}, Oberwolfach Reports, pp.\ 2879--2905, DOI \texttt{10.4171/owr/2012/47}. Supplied source URL and literature-source link, labelled ``original OWR report'': \url{https://doi.org/10.4171/owr/2012/47}. The citation's 2013 designation and the record's 2012 year are retained without resolving their relationship.
\item The supplied statement status is \texttt{corrected\_verified}. Its verification note says that the official report asks about the additional Cartan-scheme objects, rather than unidentified Nichols algebras. This is supplied verification metadata, not an independent verification here.
\item Literature checked at: \texttt{2026-08-21}. The supplied dated curation labels the problem open and reports that no complete resolution was found for the question described there in terms of unidentified diagonal Nichols braidings. This is a supplied dated assessment, not an independently verified literature conclusion. The supplied background expressly requests independent verification of time-sensitive status and cited claims.
\end{itemize}
\end{document}
